\section{Summary}\label{sec:summary}

Trapped-ion experiments with \Ybion{} deliver their \SI{369.5}{nm} cooling and detection light through
single-mode, polarisation-maintaining optical fibres. In the group, these fibres show a characteristic
failure. Directly after the connector end faces are polished, the coupling efficiency is excellent. It
then falls over days to weeks, faster at higher power, and both the input and the output end are
affected. Re-polishing restores the throughput, so the damage must sit at, or within micrometres of, the
end faces rather than in the bulk of the fibre. At a few milliwatts in a mode-field diameter of
about \SI{2.3}{\micro m}, the end faces see \SIrange{e5}{e6}{W/cm^2} of near-UV light, an
intensity at which photochemistry of the glass surface and of anything adsorbed on it becomes
important.

This project tests the hypothesis that \textbf{removing oxygen from the end faces halts the degradation},
and develops a practical way of doing so with \textbf{3D-printed enclosures} that are either purged with
inert gas or sealed with an \textbf{oxygen absorber}. The literature does not settle the question in advance. Oxygen
could drive the damage through reactions with photo-activated surface sites. In the laser-induced
contamination of UV optics in space lasers, however, oxygen is known to \emph{suppress} the growth of
absorbing deposits from organic vapours. Water vapour and airborne particles are further candidates.
The experiment is therefore designed to separate these factors. A multi-arm testbed exposes identical
fibres \emph{simultaneously} to ambient air, clean dry synthetic air, clean dry nitrogen and sealed
absorber cells, and later to controlled oxygen, humidity and contaminant levels. Throughput,
back-reflection and the atmosphere in every cell are logged continuously. At the end of each run a fixed
cleaning and re-polishing sequence, together with surface analysis, shows where the loss is and what
it consists of.

The project runs over 12 months (a 6-month variant is included). It starts with a forensic study of
fibres already degraded in the lab, builds and qualifies the testbed and cells, runs two atmosphere
campaigns, and ends with a demonstration on a working \SI{369.5}{nm} link. Success means a quantitative
answer to whether oxygen matters (with confidence intervals), and a qualified cell design that holds
$<\SI{100}{ppm}$ \Oxy{} for months. If the atmosphere is the key, the degradation should fall below
\SI{1}{\percent} per \SI{1000}{h} at \SI{3}{mW}. Part~II of this document is a practical guide for the
student: onboarding plan, core skills, working practices, pitfalls, a prioritised reading list and
draft procedures.
